\section{Ablation}
\label{sec:ablation}

Table~\ref{tab:ablation} removes or replaces one component of the proposed test at a time and reads the effect on the stated population and on the re-estimated one, in the base cells of Table~\ref{tab:main-results}, with a 2-year lengthening for the re-estimated parameter and a 2-year lengthening made one date early in the last column. The components are the conjunction with the breach flag, the exact inequality as the definition of a change, the same-date rule, and the counting statistic.

\begin{table*}[htbp]
\centering
\footnotesize
\caption{Ablation of the proposed test: each row removes or replaces one component (the conjunction with the breach flag, the exact inequality as the definition of a change, the same-date rule, the counting statistic) and reports the false-flag rate, the detection rate and the area under the ROC curve on the stated parameter and, for a 2-year lengthening, on the fitted parameter, and the detection of a 4-year lengthening made at the date before a breach begins.}
\label{tab:ablation}
\resizebox{\textwidth}{!}{%
\begin{tabular}{lccccccc}
\toprule
\multicolumn{1}{l}{} & \multicolumn{3}{c}{Stated parameter} & \multicolumn{3}{c}{Re-estimated parameter} & \multicolumn{1}{c}{Detection} \\
\cmidrule(lr){2-4}\cmidrule(lr){5-7}\cmidrule(lr){8-8}
Variant of the method & False flags (\%) & Detection (\%) & AUC & False flags (\%) & Detection (\%) & AUC & Override before the breach (\%) \\
\midrule
Proposed test as implemented: exact change, breach conjunction, same date & 34.8 & 100.0 & 0.908 & 82.2 & 100.0 & 0.579 & 81.9 \\
Without the breach conjunction (any change) & 68.8 & 100.0 & 0.863 & 92.5 & 100.0 & 0.536 & 100.0 \\
Exact change replaced by a calibrated tolerance & -- & -- & -- & 3.9 & 18.0 & 0.937 & 3.6 \\
Calibrated tolerance with a one-date look-ahead & -- & -- & -- & 4.5 & 17.2 & 0.925 & 98.6 \\
Counting replaced by an association test (Fisher) & 0.1 & 13.2 & 0.922 & 2.1 & 4.1 & 0.596 & 0.8 \\
Counting replaced by a rank test of the size of the changes (Mann--Whitney) & 2.6 & 48.0 & 0.914 & 7.2 & 31.4 & 0.796 & 1.5 \\
Counting replaced by a test for a shift in the level, whatever the breach flags & 4.0 & 5.8 & 0.796 & 5.9 & 8.4 & 0.756 & 42.0 \\
Counting replaced by a test for a shift in the level that falls on a breach date & 4.1 & 6.8 & 0.824 & 6.2 & 25.1 & 0.782 & 0.9 \\
\bottomrule
\end{tabular}}
\end{table*}


The conjunction with the breach flag is what separates the test from a monitor of parameter stability. Without it, the rule flags any history in which the parameter changed, and on the stated population the false-flag rate rises from \statedNullEngine{} to \statedNullAny{} at an unchanged detection of \statedAltAny{}; on the re-estimated population it rises from \fitNullEngine{} to \fitNullAny{} and the area under the ROC curve falls from \fitAucCEngine{} to \fitAucCAny{}. The component contributes specificity at no cost in detection, which is the reason to read the flags, and it contributes little on a re-estimated parameter, where almost every history has a change.

The exact inequality is the component with the two faces. On a stated parameter it is what gives the detection of one, since any change at all, however small, is a decision. On a re-estimated parameter it is what removes the discrimination: with the exact inequality the rule flags \fitNullEngine{} of the null banks and \fitAddCEngine{} of those with a 2-year lengthening, and replacing it by the tolerance calibrated at 5\,\% gives \fitNullTol{} and \fitAddCTol{} with an area of \fitAucCTol{}. The tolerance buys a stated size and a usable ranking at the price of the detection of small overrides, and it is the only variant that reads the size of a change at all. The same-date rule, replaced by the one-date look-ahead, does not help against an override made one date before the breach (\earlyTolLa{} against \earlyTol{} for the tolerance), because a persistent override produces one change at the date it is made and the look-ahead only widens the set of dates that can carry it by one; the cost is a small rise of the false-flag rate (\fitNullTolLa{} against \fitNullTol{}).

Replacing the count by an association test, the Fisher test, or by a rank test of the size of the changes, the Mann--Whitney test, trades detection for specificity. On the stated population they flag \statedNullFisher{} and \statedNullMw{} of the null histories and \statedAltFisher{} and \statedAltMw{} of the alternative ones, with areas of \statedAucFisher{} and \statedAucMw{}, against \statedAucEngine{} for the count: the tests rank slightly better and decide very differently, because the count has a decision at one coincidence and the tests need evidence that grows with the number of breach dates. On the re-estimated population the Mann--Whitney test has the largest area of the three published tests, \fitAucCMw{}, with a false-flag rate of \fitNullMw{} that exceeds its nominal level.

The ablation supports three readings. The conjunction with the flag and the exactness are both needed for the property that matters on a stated parameter, a flag that is certain for a change at a breach date and absent otherwise. The exactness has to be given up, for a tolerance read from a null, on a re-estimated parameter. And no component recovers an override that the data hide, as the override toward the longest tranche of Table~\ref{tab:weight-override} shows.
